\chapter{Bahasa Matematika dan Seni Membuktikan}

Di suatu kelas, seorang mahasiswa menulis di papan: ``$n^{2}+n+41$ prima untuk setiap
bilangan bulat $n \geq 0$''. Ia lalu memeriksa satu per satu. Untuk $n=0$ diperoleh
$41$; untuk $n=1$ diperoleh $43$; untuk $n=2$ diperoleh $47$. Ia melanjutkan sampai
$n=39$, mendapat $1601$, dan $1601$ memang prima. Empat puluh nilai, empat puluh
kegagalan yang tidak terjadi. Di $n=40$ klaim itu patah:
$40^{2}+40+41 = 1681 = 41^{2}$. Satu baris tambahan membatalkan seluruh halaman
sebelumnya.

Kejadian itu bukan kecelakaan aritmetika, melainkan kesalahan tata bahasa. ``Untuk setiap
$n \geq 0$'' berbicara tentang tak berhingga banyak bilangan, sedangkan daftar hanya
menyentuh berhingga banyak bilangan; tidak ada aturan logika yang memindahkan kebenaran
dari yang berhingga ke yang tak berhingga. Bab ini karena itu menawarkan cara berbicara
sehingga pemeriksaan tidak perlu.

Kebutuhan itu bukan kemewahan. Di Bab~29 kita akan menyimpulkan bahwa sebuah
pengklasifikasi dapat dikelilingi dengan menggeser satu masukan sejauh $\varepsilon$,
dan kesimpulan itu berlaku untuk semua masukan. Di Bab~27 laporan akurasi yang bagus pada
ribuan data uji tidak membuktikan apa pun tentang data esok hari. Bab ini menyiapkan
perangkat untuk keduanya.

\begin{untukanak}
Di tepi sebuah danau hiduplah sekawanan angsa. Seorang anak datang setiap pagi selama
hampir setahun dan menghitung angsa yang dilihatnya: putih, semuanya, selama dua ratus
tujuh puluh pagi. Ia pun menulis di buku catatannya: semua angsa di danau ini putih.

Kakeknya bertanya berapa angsa yang belum pernah dilihat anak itu. Anak itu mengaku
tidak tahu, dan kakeknya berkata bahwa catatan itu bercerita tentang angsa yang sudah
dilihat, bukan tentang semua angsa di danau.

Cerita ini bukan tentang angsa, melainkan tentang beda antara ``sudah kuperiksa
semuanya'' dan ``berlaku untuk semuanya''. Yang pertama hanya sekuat ingatan kita. Yang kedua tidak bergantung pada siapa yang memeriksa. Angsa hitam
cukup satu untuk meruntuhkan catatan setebal apa pun, sedangkan seribu angsa putih
tidak pernah cukup untuk menegakkannya.

Angsa hitam pada bab ini bernama empat puluh. Pada bilangan itu rumus yang tampak rapi
mendadak menghasilkan bilangan yang bisa dipecah. Di bab ini kita akan tahu mengapa satu
contoh meruntuhkan klaim, mengapa banyak contoh tidak menegakkannya, dan bagaimana celah
itu ditutup tanpa memeriksa angsa terakhir.
\end{untukanak}

\section{Pernyataan, penghubung, dan tabel kebenaran}
\tingkat{dasar}

Kalimat baru bisa dinilai setelah kita tahu apa yang dijanjikannya. ``Tujuh adalah
bilangan prima'' menjanjikan sesuatu yang bisa diperiksa, dan janji itu benar. ``Tutup
jendelanya'' tidak menjanjikan apa pun. ``$x+1=3$'' belum bisa dinilai, karena nilainya
bergantung pada $x$. Ketiganya berbentuk sama di mata ejaan, tetapi hanya yang pertama
termasuk bahan matematika.

\begin{definisi}[Pernyataan]
\index{pernyataan}
\emph{Pernyataan} adalah kalimat yang bernilai benar atau bernilai salah, tetapi tidak
keduanya sekaligus.
\end{definisi}

Pernyataan tunggal menjadi berguna setelah digabung. Setiap penghubung punya satu arti
yang ditetapkan, bukan disimpulkan dari kebiasaan berbahasa.

\kamus{penghubung logika}{logical connective}{seperti tombol pada mesin hitung: dua
lampu masuk, satu lampu keluar, dan hubungan antarlampu ditetapkan pabrik, bukan
ditebak pemakainya.}

Dengan gambar tombol itu di kepala, kita tetapkan arti kelima penghubung secara resmi.

\begin{definisi}[Penghubung logika]
Misalkan $P$ dan $Q$ pernyataan.

Lima penghubung berikut kita pakai.
\begin{itemize}
  \item $\neg P$ (\emph{negasi}, \emph{not}): benar tepat ketika $P$ salah.
  \item $P \wedge Q$ (\emph{konjungsi}, \emph{and}): benar tepat ketika keduanya benar.
  \item $P \vee Q$ (\emph{disjungsi}, \emph{or}): benar tepat ketika paling sedikit
        satu benar; kata ``atau'' di sini inklusif.
  \item $P \Rightarrow Q$ (\emph{implikasi}, \emph{implication}): salah tepat ketika
        $P$ benar dan $Q$ salah; benar pada tiga kemungkinan lainnya.
  \item $P \Leftrightarrow Q$ (\emph{bikondisional}, \emph{iff}): benar tepat ketika
        $P$ dan $Q$ bernilai sama.
\end{itemize}
\end{definisi}

Seluruh isi definisi itu dapat dipadatkan menjadi satu tabel: barisnya semua kombinasi
nilai, kolomnya nilai setiap penghubung.

{\centering
\begin{tabular}{cc|ccccc}
\hline
$P$ & $Q$ & $\neg P$ & $P \wedge Q$ & $P \vee Q$ & $P \Rightarrow Q$ & $P \Leftrightarrow Q$ \\
\hline
B & B & S & B & B & B & B \\
B & S & S & S & B & S & S \\
S & B & B & S & B & B & S \\
S & S & B & S & S & B & B \\
\hline
\end{tabular}\par}

\begin{figure}[htbp]
\centering
\begin{tikzpicture}[font=\small]
  \fill[layerjade!16] (0,0) rectangle (3,1.8);
  \fill[accent!16] (3,0) rectangle (6,1.8);
  \fill[layerjade!16] (0,-1.8) rectangle (3,0);
  \fill[layerjade!16] (3,-1.8) rectangle (6,0);
  \draw[layergrey!70,thick] (0,-1.8) grid[xstep=3,ystep=1.8] (6,1.8);
  \draw[accent,very thick] (3,0) rectangle (6,1.8);
  \node at (1.5,0.9) {B};
  \node at (4.5,0.9) {S};
  \node at (1.5,-0.9) {B};
  \node at (4.5,-0.9) {B};
  \node[above,font=\footnotesize] at (1.5,1.8) {$Q$ benar};
  \node[above,font=\footnotesize] at (4.5,1.8) {$Q$ salah};
  \node[left,font=\footnotesize] at (0,0.9) {$P$ benar};
  \node[left,font=\footnotesize] at (0,-0.9) {$P$ salah};
  \node[below,font=\footnotesize,text=accent] at (4.5,-1.8) {satu-satunya sel yang salah};
\end{tikzpicture}
\caption{Kisi kebenaran $P \Rightarrow Q$: baris berisi nilai $P$, kolom berisi nilai $Q$.
Tiga sel bernilai B dan satu sel bernilai S, yaitu ketika $P$ benar sementara $Q$ salah.
Siapa pun yang ingin menggugurkan implikasi cukup menunjuk sel yang dibingkai tebal.}
\label{fig:kisi-implikasi}
\end{figure}

Baris ketiga dan keempat adalah yang paling sering mengejutkan. Pada baris keempat,
$P$ salah dan $Q$ salah, dan implikasinya dinilai benar. Pada baris ketiga, $P$ salah
dan $Q$ benar, dan implikasinya juga benar. Jadi satu-satunya cara menggugurkan
$P \Rightarrow Q$ adalah menunjukkan $P$ benar sementara $Q$ salah; kisi pada
Gambar~\ref{fig:kisi-implikasi} memuat sel itu sebagai satu-satunya sel yang salah.

Ketetapan itu bukan pilihan gaya. Perhatikan ``jika $n$ habis dibagi $4$, maka $n$
genap'', yang harus berlaku untuk setiap bilangan bulat $n$, termasuk $n=3$. Di situ
antesedennya salah; kalau implikasi dinilai salah, pernyataan yang jelas benar tadi akan
gugur karena satu bilangan yang tidak relevan. Implikasi mencatat hubungan nilai, bukan
sebab-akibat.

\kamus{implikasi}{implication}{seperti janji ``kalau hujan, jalan basah'': janji itu
hanya dilanggar bila hujan turun dan jalan tetap kering; bila tidak hujan, si pemilik
janji belum berbuat salah.}

Dua pernyataan majemuk bisa bernilai sama pada semua kemungkinan; kesamaan semacam itu
punya nama sendiri.

\begin{definisi}[Tautologi dan kesetaraan]
\index{tautologi}
Pernyataan majemuk disebut \emph{tautologi} jika bernilai benar pada setiap penetapan
nilai pernyataan-pernyataan penyusunnya. Dua pernyataan majemuk $A$ dan $B$ disebut
\emph{setara secara logis}, ditulis $A \equiv B$, jika $A \Leftrightarrow B$ tautologi.
\end{definisi}

Untuk dua pernyataan penyusun ada $2^{2}=4$ penetapan nilai; untuk tiga penyusun ada
$2^{3}=8$. Kesetaraan yang hanya melibatkan dua atau tiga pernyataan karena itu dapat
dibuktikan dengan memeriksa seluruh baris, dan pemeriksaan itu lengkap. Inilah satu-satunya tempat dalam bab ini di mana
pemeriksaan habis-habisan diterima sebagai bukti, sebab himpunan kasusnya berhingga dan
sudah dihitung.

\begin{proposisi}
Untuk setiap pernyataan $P$ dan $Q$ berlaku $P \Rightarrow Q \equiv \neg P \vee Q$.
\end{proposisi}
\par\smallskip\noindent\textit{Bukti.} Kita bandingkan kolom $P \Rightarrow Q$ dan kolom $\neg P \vee Q$ pada keempat baris
tabel kebenaran. Pada baris pertama keduanya benar. Pada baris kedua $\neg P$ salah dan
$Q$ salah, sehingga keduanya salah. Pada baris ketiga dan keempat $\neg P$ benar,
sehingga keduanya benar. Keempat baris memberi nilai yang sama, jadi kedua pernyataan
setara.
\hfill$\blacksquare${\parfillskip=0pt\par}

Dua kesetaraan turunan dipakai sepanjang buku ini. Yang pertama membolehkan kita membalik
arah pembicaraan tanpa mengubah nilai.

\begin{akibat}[Kontraposisi]
\index{kontraposisi}
Untuk setiap pernyataan $P$ dan $Q$ berlaku
$P \Rightarrow Q \equiv \neg Q \Rightarrow \neg P$.
\end{akibat}
\par\smallskip\noindent\textit{Bukti.} Menurut proposisi sebelumnya, kedua ruas sama dengan $\neg P \vee Q$ dan
$\neg(\neg Q) \vee \neg P$. Karena $\neg(\neg Q) \equiv Q$ dan disjungsi tidak
memperhatikan urutan, kedua bentuk itu identik.
\hfill$\blacksquare${\parfillskip=0pt\par}

Yang \emph{tidak} setara adalah konversnya, yaitu $Q \Rightarrow P$. Pada baris ketiga
tabel, $Q \Rightarrow P$ benar sementara $P \Rightarrow Q$ salah. Menukar anteseden
dengan konsekuen adalah kesalahan tersering, dan sering tidak terlihat karena kalimatnya
tetap enak dibaca.

\begin{proposisi}
Untuk setiap pernyataan $P$ dan $Q$ berlaku
$\neg(P \Rightarrow Q) \equiv P \wedge \neg Q$.
\end{proposisi}
\par\smallskip\noindent\textit{Bukti.} Menurut definisi negasi dan proposisi pertama,
$\neg(P \Rightarrow Q) \equiv \neg(\neg P \vee Q)$. Hukum De Morgan yang dibuktikan di
Bab~2 memberi $\neg(\neg P \vee Q) \equiv \neg(\neg P) \wedge \neg Q$, dan
$\neg(\neg P) \equiv P$. Jadi $\neg(P \Rightarrow Q) \equiv P \wedge \neg Q$.
Pembaca yang belum menerima hukum De Morgan dapat memeriksa keempat baris tabel secara
langsung.
\hfill$\blacksquare${\parfillskip=0pt\par}

\begin{tangan}
Sediakan delapan kartu kecil: empat bertuliskan B, empat bertuliskan S. Susun empat baris,
setiap baris satu kartu untuk $P$ dan satu untuk $Q$: B-B, B-S, S-B, S-S. Di kanan setiap
baris letakkan kartu ketiga yang Anda tulis sendiri: B kalau ``jika $P$ maka $Q$'' Anda
nilai benar, S kalau salah.

Kerjakan tanpa melihat tabel, lalu bandingkan. Hanya satu baris mendapat kartu S, yaitu
baris yang $P$-nya B dan $Q$-nya S. Geser kartu itu ke atas dan ke bawah; ia tetap
satu-satunya. Sekarang tukar kartu P dan Q pada baris B-S sehingga menjadi S-B, lalu nilai
ulang: kartunya berubah dari S menjadi B. Perpindahan satu kartu itulah konvers.
\end{tangan}

\section{Kuantor dan cara menegasikannya}
\tingkat{dasar}

Sebagian besar kalimat matematika berbicara tentang banyak objek, dan banyaknya itu
ditandai kuantor. Ada dua kuantor, dan keduanya membaca urutan sebagai bagian dari arti.

\kamus{kuantor}{quantifier}{seperti dua jenis pengumuman di sekolah: ``semua murid
wajib membawa buku'' dan ``ada murid yang membawa buku''; keduanya terdengar mirip
tetapi menuntut bukti yang sangat berbeda.}

Kita tuliskan artinya secara resmi.

\begin{definisi}[Kuantor universal dan eksistensial]
\index{kuantor}
Misalkan $P(x)$ pernyataan yang bergantung pada $x$ di dalam himpunan $A$.
\begin{itemize}
  \item $(\forall x \in A)\,P(x)$ berarti $P(x)$ benar untuk setiap $x \in A$.
  \item $(\exists x \in A)\,P(x)$ berarti terdapat paling sedikit satu $x \in A$
        dengan $P(x)$ benar.
\end{itemize}
Pernyataan pertama salah jika ada satu $x \in A$ dengan $P(x)$ salah; satu objek
semacam itu disebut \emph{contoh penyangkal} (\emph{counterexample}).
\end{definisi}

Kuantor eksistensial menuntut paling sedikit satu, bukan tepat satu; menambahkan tuntutan
yang tidak tercantum dalam definisi adalah sumber kesalahan tersering.

Urutan dua kuantor yang berbeda jenis tidak boleh ditukar. Di dalam $\mathbb{R}$,
pernyataan
\[
  (\forall x)(\exists y)\; y > x
\]
benar, sedangkan
\[
  (\exists y)(\forall x)\; y > x
\]
salah. Bukti keduanya pendek.

\begin{proposisi}
Di dalam $\mathbb{R}$ berlaku $(\forall x)(\exists y)\, y > x$, tetapi tidak berlaku
$(\exists y)(\forall x)\, y > x$.
\end{proposisi}
\par\smallskip\noindent\textit{Bukti.} Untuk bagian pertama, ambil $x \in \mathbb{R}$
sebarang dan pilih $y = x+1$. Maka $y > x$, sehingga $y$ menjadi saksi bagi $x$ itu.
Karena $x$ sebarang, setiap $x$ punya saksi; saksinya boleh berbeda untuk $x$ yang
berbeda.

Untuk bagian kedua, andaikan ada $y_{0} \in \mathbb{R}$ dengan $y_{0} > x$ untuk
setiap $x \in \mathbb{R}$. Ambil $x = y_{0}+1$, yang juga bilangan real. Maka
$y_{0} > y_{0}+1$, yaitu $0 > 1$, suatu hal yang tidak mungkin. Jadi tidak ada $y_{0}$
semacam itu.
\hfill$\blacksquare${\parfillskip=0pt\par}

\index{negasi}Negasi bekerja dari luar ke dalam: setiap kuantor berganti jenis, dan negasi berjalan masuk
sampai menempel pada pernyataan terdalam. Untuk pernyataan majemuk, hukum De Morgan
memindahkan negasi melewati konjungsi dan disjungsi.

{\centering
\begin{tabular}{ll}
\hline
Pernyataan & Negasinya \\
\hline
$(\forall x \in A)\,P(x)$ & $(\exists x \in A)\,\neg P(x)$ \\
$(\exists x \in A)\,P(x)$ & $(\forall x \in A)\,\neg P(x)$ \\
$P \wedge Q$ & $\neg P \vee \neg Q$ \\
$P \vee Q$ & $\neg P \wedge \neg Q$ \\
$P \Rightarrow Q$ & $P \wedge \neg Q$ \\
\hline
\end{tabular}\par}

Baris terakhir itu yang paling sering dilanggar. Menyangkal ``jika hujan maka jalan
basah'' bukan ``jika hujan maka jalan kering'', melainkan ``hujan dan jalan tidak basah''.
Kesalahan yang sama muncul saat orang menyangkal definisi: kata kerja di dalamnya
dinegasikan, kuantornya dibiarkan utuh.

\begin{proposisi}
Untuk setiap pernyataan $P(x)$ dan $Q(x)$ berlaku
\[
  \neg(\forall x)\bigl(P(x) \Rightarrow Q(x)\bigr)
  \;\equiv\; (\exists x)\bigl(P(x) \wedge \neg Q(x)\bigr).
\]
\end{proposisi}
\par\smallskip\noindent\textit{Bukti.} Kuantor universal dinegasikan menjadi kuantor
eksistensial atas negasi badan:
\[
  \neg(\forall x)R(x) \;\equiv\; (\exists x)\neg R(x).
\]
Badan $R(x)$ di sini adalah $P(x) \Rightarrow Q(x)$, yang negasinya $P(x) \wedge
\neg Q(x)$ menurut proposisi bagian sebelumnya. Jadi negasinya berbentuk
$(\exists x)(P(x) \wedge \neg Q(x))$.
\hfill$\blacksquare${\parfillskip=0pt\par}

Bentuk itu memberi tugas yang jelas ketika hendak menolak pernyataan berkuantor: tunjukkan
satu objek yang memenuhi anteseden dan gagal pada konsekuen.

Dua negasi berikut dipakai berulang kali; yang pertama menjadi bahan Bab~4.

\begin{contoh}[Negasi definisi limit]
Definisi limit berbunyi: barisan $(a_{n})$ menuju $L$ jika
\[
  (\forall \varepsilon > 0)(\exists N \in \mathbb{N})(\forall n \geq N)\;
  |a_{n} - L| < \varepsilon .
\]
Menegasikan dari luar ke dalam memberi
\[
  (\exists \varepsilon > 0)(\forall N \in \mathbb{N})(\exists n \geq N)\;
  |a_{n} - L| \geq \varepsilon .
\]
Bacaan kalimat kedua: ada satu toleransi $\varepsilon$ yang tidak bisa diperbaiki oleh
$N$ mana pun, karena selalu ada suku setelah $N$ yang menyimpang paling sedikit sejauh
$\varepsilon$. Untuk membuktikan bahwa $(a_{n})$ tidak menuju $L$, kita harus menemukan
$\varepsilon$ itu, lalu untuk setiap $N$ menemukan $n \geq N$ yang melanggar; $n$ boleh
bergantung pada $N$.
\end{contoh}

Negasi kedua menyangkut keberadaan bilangan prima. Euklides membuktikan bahwa untuk
setiap bilangan asli $n$ ada bilangan prima yang lebih besar daripada $n$:
\[
  (\forall n \in \mathbb{N})(\exists p)\bigl(p \text{ prima} \wedge p > n\bigr).
\]
Negasinya adalah
\[
  (\exists n \in \mathbb{N})(\forall p)\bigl(p \text{ prima} \Rightarrow p \leq n\bigr),
\]
yaitu: ada satu batas $n$ sehingga semua bilangan prima tidak melebihi $n$. Menyangkal
klaim Euklides menuntut batas semacam itu, bukan sekadar gagal menemukan prima besar.

\begin{untukanak}
Bayangkan sebuah kelas dengan tiga puluh murid. Guru menulis di papan: ``semua murid
membawa pensil.'' Untuk membuktikan kalimat itu benar, guru harus memeriksa ketiga puluh meja; untuk
membuktikannya salah, cukup satu murid yang tidak membawa pensil.

Sekarang guru menulis kalimat lain: ``ada murid yang membawa pensil.'' Kali ini
keadaannya terbalik. Satu murid yang membawa pensil sudah cukup untuk membuktikan
kalimat itu. Tetapi untuk membuktikannya salah, guru harus memeriksa ketiga puluh meja
dan memastikan tidak ada satu pun pensil.

Dua kalimat itu hanya bertukar kata di depan, tetapi pekerjaan memeriksanya bertukar
seluruhnya. Ahli matematika berhati-hati dengan ``semua'' dan ``ada'' karena keduanya
menghitung berapa banyak yang harus diperiksa.

Ada satu jebakan lagi. Kalau ada murid yang membawa dua pensil, apakah ``ada murid yang
membawa pensil'' masih benar? Tentu benar: kalimat itu menuntut paling sedikit satu,
bukan tepat satu. Di bab ini kita akan tahu mengapa tuntutan yang tidak tertulis tidak
boleh ditambahkan sendiri.
\end{untukanak}

\section{Membaca definisi secara harfiah}
\tingkat{menengah}

Definisi matematika bukan penjelasan singkat tentang sebuah kata; ia kontrak yang
berlaku dua arah. Kalau sebuah bilangan disebut ganjil apabila berbentuk $2k+1$ untuk
suatu bilangan bulat $k$, kita terikat pada dua hal: setiap bilangan berbentuk $2k+1$
harus disebut ganjil, dan setiap bilangan yang disebut ganjil harus berbentuk $2k+1$.

Membaca definisi berarti membaca satu kalimat berbentuk ``$x$ disebut $B$ jika dan hanya
jika $C(x)$'', walaupun kata ``hanya jika'' jarang ditulis; frasa ``disebut'' sudah
mengandung kedua arah. Arah pertama memberi izin memakai nama $B$, arah kedua memberi izin
memakai syarat $C(x)$, dan membalik salah satunya merusak seluruh argumen.

\kamus{syarat perlu}{necessary condition}{seperti syarat masuk gedung: tanpa kartu
anggota kita tidak boleh masuk, tetapi punya kartu anggota belum berarti kita sudah
duduk di dalam.}

Bandingkan dua pernyataan tentang bilangan bulat $n$: ``jika $6$ membagi $n$ maka $3$
membagi $n$'' benar, sedangkan ``jika $3$ membagi $n$ maka $6$ membagi $n$'' salah,
dengan contoh penyangkal $n=3$. Jadi ``$3$ membagi $n$'' syarat perlu bagi ``$6$ membagi
$n$'', tetapi bukan syarat cukup. Definisi yang baik menyebut keduanya, sebab hanya
dengan begitu nama dan syaratnya dapat ditukar.

\begin{bacakeras}
Bacakan sebagai dialog dua suara, bergantian, dengan satu tarikan napas di setiap titik.

Suara pertama: ``Kalau hujan, jalan basah.''
Suara kedua: ``Jalan basah. Berarti hujan.''
Suara pertama: ``Belum tentu. Jalan juga basah kalau disiram.''
Suara kedua: ``Tetap basah. Tetap hujan.''
Suara pertama: ``Yang kujanjikan hanya satu arah. Hujan dulu, baru basah. Arah balik tidak pernah kujanjikan.''

Ulangi dialog itu sekali lagi, lebih cepat, sampai telinga menangkap bedanya. Kalimat
pertama menaruh hujan di depan; kalimat kedua memindahkan hujan ke belakang. Bunyinya
berputar, janjinya tidak.

Baris yang harus diulang tiga kali dengan suara keras: ``jika hujan maka basah, bukan jika
basah maka hujan.''
\end{bacakeras}

Sekarang perhatikan definisi yang paling banyak dipakai di buku ini.

\begin{definisi}[Bilangan prima]
\index{bilangan prima}
Bilangan bulat $p$ disebut \emph{prima} jika $p \geq 2$ dan satu-satunya pembagi
positif $p$ adalah $1$ dan $p$ sendiri.
\end{definisi}

Kalimat sehari-harinya: bilangan prima lebih besar daripada satu dan tidak dapat dipecah
menjadi hasil kali dua bilangan yang lebih kecil; kedua klausanya diuji terpisah.

\begin{proposisi}
Bilangan $1$ bukan bilangan prima.
\end{proposisi}
\par\smallskip\noindent\textit{Bukti.} Definisi menuntut $p \geq 2$. Untuk $p = 1$
syarat itu gagal, sehingga $1$ tidak memenuhi definisi, apa pun yang dapat kita
katakan tentang pembaginya. Perhatikan bahwa $1$ sebenarnya lolos pada klausa kedua:
satu-satunya pembagi positif $1$ memang $1$ dan $1$ sendiri. Jadi yang menyingkirkannya
adalah klausa pertama, bukan klausa kedua.
\hfill$\blacksquare${\parfillskip=0pt\par}

Kejadian semacam ini menjelaskan mengapa definisi perlu dibaca sampai klausa terakhir.
Pembaca yang melewatkan klausa pertama akan menyimpulkan bahwa $1$ prima, dan
ketunggalan faktorisasi $12 = 2^{2}\cdot 3$ di Bab~28 akan hilang, karena $12$ juga
boleh ditulis $1 \cdot 2^{2} \cdot 3$.

Definisi juga menentukan cara menyangkal keanggotaan. Untuk menunjukkan bahwa sebuah
bilangan bukan prima, kita harus menemukan pembagi positif yang bukan $1$ dan bukan
bilangan itu sendiri. Berikut satu contoh yang sering lolos dari pemeriksaan cepat.

\begin{proposisi}
Bilangan $91$ bukan bilangan prima. Lebih umum, pernyataan ``setiap bilangan ganjil
yang tidak habis dibagi $3$ dan tidak habis dibagi $5$ adalah prima'' salah.
\end{proposisi}
\par\smallskip\noindent\textit{Bukti.} Faktanya $91 = 7 \cdot 13$, dan kedua faktor itu
berada di antara $1$ dan $91$, sehingga $91$ punya pembagi positif selain $1$ dan
dirinya. Bilangan $91$ ganjil, tidak habis dibagi $3$ karena $9+1 = 10$ bukan kelipatan
$3$, dan tidak habis dibagi $5$ karena angka terakhirnya bukan $0$ atau $5$. Jadi $91$
adalah contoh penyangkal bagi pernyataan umum tadi. Sebagai pembanding, $97$ memang prima; di antara bilangan bulat $2$ dan $99$ terdapat
tepat $25$ bilangan prima, hasil hitungan program pada bagian lab.
\hfill$\blacksquare${\parfillskip=0pt\par}

\begin{definisi}[Kuadrat sempurna]
Bilangan bulat $n$ disebut \emph{kuadrat sempurna} jika terdapat bilangan bulat $k$
dengan $n = k^{2}$.
\end{definisi}

Menurut definisi ini, $0$ adalah kuadrat sempurna, karena $0 = 0^{2}$, dan $1$ juga,
karena $1 = 1^{2}$. Sementara itu $2$ bukan kuadrat sempurna, dan pernyataan itu butuh
bukti, bukan pengakuan bahwa kita tidak menemukan $k$.

\begin{proposisi}
Bilangan $2$ bukan kuadrat sempurna.
\end{proposisi}
\par\smallskip\noindent\textit{Bukti.} Andaikan $2 = k^{2}$ untuk suatu $k \in
\mathbb{Z}$. Karena $k^{2} = 2 > 0$, berlaku $k \neq 0$. Jika $|k| = 1$ maka
$k^{2} = 1 \neq 2$. Jika $|k| \geq 2$ maka $k^{2} \geq 4 > 2$. Karena $|k|$ bilangan
bulat positif, hanya dua kemungkinan itu yang tersisa, dan keduanya gagal. Jadi tidak
ada $k$ semacam itu.
\hfill$\blacksquare${\parfillskip=0pt\par}

\index{contoh penyangkal}Kebiasaan yang layak dipelihara: setiap kali bertemu
definisi, ajukan tiga pertanyaan. Apa saja klausanya? Objek apa yang memenuhi
seluruhnya? Objek apa yang sering dianggap memenuhi padahal gagal pada satu klausa?
Pertanyaan ketiga menghasilkan contoh penyangkal, dan contoh penyangkal tidak dapat
dibantah maupun diperdebatkan.

\section{Empat pola bukti}
\tingkat{menengah}

Bukti adalah pengangkutan kebenaran: dari pernyataan yang sudah kita terima menuju
pernyataan yang ingin kita terima, lewat langkah yang masing-masing dapat diperiksa.
Empat pola di bawah ini menutup sebagian besar kebutuhan, dan
Gambar~\ref{fig:peta-tiga-bukti} memetakan tiga jalur pertamanya. Pola kelima yang lazim,
pembagian kasus, hanyalah pola pertama yang diterapkan pada kasus-kasus yang saling
meniadakan.

\begin{figure}[htbp]
\centering
\begin{tikzpicture}[font=\small,>=Stealth,
  jalur/.style={draw,rounded corners=2pt,align=center,inner sep=5pt,text width=4.6cm},
  ujung/.style={draw,rounded corners=2pt,align=center,inner sep=4pt,text width=2.2cm}]
  \node[jalur,fill=layerorange!12] (klaim) at (0,0) {klaim: $P \Rightarrow Q$};
  \node[jalur,fill=layerteal!12] (lgs) at (6.1,2.4)
    {\textbf{langsung}\\ uraikan $P$, hitung, muncullah $Q$};
  \node[jalur,fill=layerindigo!10] (ktp) at (6.1,0)
    {\textbf{kontraposisi}\\ buktikan $\neg Q \Rightarrow \neg P$; nilainya sama};
  \node[jalur,fill=layerplum!10] (ktr) at (6.1,-2.4)
    {\textbf{kontradiksi}\\ andaikan $P$ benar dan $Q$ salah, turunkan pertentangan};
  \node[ujung,fill=layergrey!10] (selesai) at (11.6,0)
    {tidak ada\\penetapan yang\\menggugurkan};
  \draw[->,thick] (klaim.east) -- (lgs.west);
  \draw[->,thick] (klaim.east) -- (ktp.west);
  \draw[->,thick] (klaim.east) -- (ktr.west);
  \draw[->,thick] (lgs.east) -- (selesai.north west);
  \draw[->,thick] (ktp.east) -- (selesai.west);
  \draw[->,thick] (ktr.east) -- (selesai.south west);
\end{tikzpicture}
\caption{Peta jalan tiga bentuk bukti implikasi. Semua jalur berangkat dari klaim yang sama
dan berakhir pada kenyataan bahwa tidak ada penetapan nilai yang membuat $P$ benar sementara
$Q$ salah. Bukti kontraposisi sah karena $P \Rightarrow Q \equiv \neg Q \Rightarrow \neg P$.
Bentuk keempat, bukti konstruktif dengan saksi, tidak masuk peta ini karena ia membuktikan
pernyataan berbentuk $\exists$, bukan implikasi.}
\label{fig:peta-tiga-bukti}
\end{figure}

\paragraph{Pola pertama: bukti langsung.} Kita berangkat dari hipotesis, menguraikannya
menurut definisi, lalu menghitung sampai kesimpulan muncul. Contoh berikut dipilih
karena kesimpulannya lebih kuat daripada yang biasanya dibuktikan orang.

\begin{teorema}
Untuk setiap bilangan bulat ganjil $n$, bilangan $n^{2}-1$ habis dibagi $8$.
\end{teorema}
\par\smallskip\noindent\textit{Bukti.} Karena $n$ ganjil, $n = 2k+1$ untuk suatu
$k \in \mathbb{Z}$. Maka
\[
  n^{2} - 1 = (2k+1)^{2} - 1 = 4k^{2} + 4k = 4k(k+1).
\]
Di antara dua bilangan bulat berurutan $k$ dan $k+1$, tepat satu di antaranya genap,
sehingga $k(k+1) = 2m$ untuk suatu $m \in \mathbb{Z}$. Substitusi memberi
$n^{2}-1 = 4 \cdot 2m = 8m$, yang berarti $n^{2}-1$ habis dibagi $8$.
\hfill$\blacksquare${\parfillskip=0pt\par}

Akibat langsungnya, kuadrat bilangan ganjil selalu bersisa $1$ bila dibagi $8$. Untuk
$n = 1,3,5,7,9,11$ berturut-turut diperoleh $n^{2}-1 = 0, 8, 24, 48, 80, 120$, dan
keenamnya habis dibagi $8$. Pemeriksaan itu bukan bukti, tetapi ia memastikan bahwa
bukti di atas tidak salah ketik.

\begin{figure}[htbp]
\centering
\begin{tikzpicture}[font=\small,>=Stealth,
  tahap/.style={draw,rounded corners=2pt,align=left,inner sep=5pt,text width=6.6cm},
  nama/.style={font=\footnotesize\itshape,text=layergrey,align=right,text width=2.4cm}]
  \node[tahap,fill=layerteal!10] (asumsi) at (0,0)
    {$n$ ganjil, jadi $n = 2k+1$ untuk suatu $k \in \mathbb{Z}$};
  \node[nama] at (-4.9,0) {asumsi (diberikan)};
  \node[tahap,fill=layerorange!10] (langkah) at (0,-2.1)
    {$n^{2}-1 = 4k^{2}+4k = 4k(k+1) = 8m$};
  \node[nama] at (-4.9,-2.1) {langkah (dihitung)};
  \node[tahap,fill=layerjade!10] (kesimpulan) at (0,-4.2)
    {$8$ membagi $n^{2}-1$};
  \node[nama] at (-4.9,-4.2) {kesimpulan (diminta)};
  \draw[->,thick] (asumsi.south) -- (langkah.north);
  \draw[->,thick] (langkah.south) -- (kesimpulan.north);
\end{tikzpicture}
\caption{Anatomi satu bukti langsung untuk pernyataan ``jika $n$ ganjil maka $8$ membagi
$n^{2}-1$''. Kotak atas satu-satunya yang boleh dipakai tanpa alasan; kotak tengah memuat
alasannya, yaitu definisi bilangan ganjil dan aljabar; kotak bawah adalah yang diminta.
Panahnya hanya boleh mengalir ke bawah.}
\label{fig:anatomi-bukti}
\end{figure}

Gambar~\ref{fig:anatomi-bukti} memberi nama ketiga bagian bukti itu: asumsi, langkah,
kesimpulan.

\paragraph{Pola kedua: bukti kontraposisi.} Kita membuktikan $\neg Q \Rightarrow \neg P$
sebagai ganti $P \Rightarrow Q$. Menurut akibat pada bagian pertama, keduanya setara,
sehingga sah. Pola ini dipilih ketika negasi kesimpulan lebih mudah dipegang daripada
hipotesisnya.

\begin{proposisi}
Untuk setiap bilangan bulat $n$, jika $5n+3$ genap maka $n$ ganjil.
\end{proposisi}
\par\smallskip\noindent\textit{Bukti.} Kontraposisinya berbunyi: jika $n$ genap maka
$5n+3$ ganjil. Misalkan $n = 2k$. Maka $5n+3 = 10k+3 = 2(5k+1)+1$, yang berbentuk dua
kali bilangan bulat ditambah satu, jadi ganjil. Karena kontraposisi terbukti,
implikasi aslinya terbukti.
\hfill$\blacksquare${\parfillskip=0pt\par}

Membuktikan langsung akan menuntut pembagian oleh $5$ dari $5n+3 = 2m$; bukan mustahil,
tetapi lebih panjang. Pola kontraposisi dipilih karena menghemat.

Sebelum pola ketiga, kita butuh satu lemma tentang paritas. Lemma ini akan dipakai dua
kali: sekali di sini, sekali pada bukti ketakrasionalan akar dua.

\begin{lemma}
\index{lemma paritas}
Untuk setiap bilangan bulat $n$, jika $n^{2}$ genap maka $n$ genap.
\end{lemma}
\par\smallskip\noindent\textit{Bukti.} Kita buktikan kontraposisinya: jika $n$ ganjil
maka $n^{2}$ ganjil. Menurut teorema pertama bagian ini, untuk $n$ ganjil berlaku
$n^{2} = 8m+1$ untuk suatu $m$, sehingga $n^{2}$ bersisa $1$ bila dibagi $2$, yaitu
ganjil.
\hfill$\blacksquare${\parfillskip=0pt\par}

\paragraph{Pola ketiga: bukti kontradiksi.} Kita mengandaikan kesimpulan salah, lalu menurunkan dua pernyataan yang tidak dapat benar
bersama. Pola ini dipakai ketika objek yang disangkal tidak punya bentuk yang nyaman
dihitung, misalnya ketika yang disangkal adalah keberadaan sesuatu.

\begin{teorema}
\index{akar dua}
Tidak ada bilangan rasional $r$ dengan $r^{2} = 2$.
\end{teorema}
\par\smallskip\noindent\textit{Bukti.} Andaikan ada $r \in \mathbb{Q}$ dengan
$r^{2} = 2$. Setiap bilangan rasional dapat ditulis sebagai $p/q$ dengan $p,q$ bilangan
bulat, $q \neq 0$, dan $p,q$ tidak memiliki faktor persekutuan selain $1$; kita
membuang faktor persekutuan itu dengan membagi keduanya, sehingga bentuk tanpa faktor
persekutuan selalu dapat dicapai.

Dari $p^{2}/q^{2} = 2$ diperoleh $p^{2} = 2q^{2}$, sehingga $p^{2}$ genap. Menurut
lemma di atas, $p$ genap; tulis $p = 2m$ dengan $m \in \mathbb{Z}$. Substitusi memberi
$4m^{2} = 2q^{2}$, yaitu $q^{2} = 2m^{2}$, sehingga $q^{2}$ genap dan, sekali lagi
menurut lemma, $q$ genap. Kini $p$ dan $q$ keduanya genap, jadi keduanya habis dibagi
$2$, bertentangan dengan pilihan $p/q$ tanpa faktor persekutuan. Pengandaian salah,
maka tidak ada $r$ rasional dengan $r^{2} = 2$.
\hfill$\blacksquare${\parfillskip=0pt\par}

Kontradiksinya bukan antara $r^{2}=2$ dan aritmetika biasa, melainkan antara ``$p,q$
keduanya genap'' dan janji awal ``$p,q$ tanpa faktor persekutuan''. Karena itu langkah
membuang faktor persekutuan bukan hiasan; ia bagian dari hipotesis yang nanti dilanggar.

\paragraph{Pola keempat: bukti konstruktif dengan saksi.} Untuk membuktikan pernyataan
berbentuk $\exists$, kita menunjuk satu objek dan memeriksa bahwa objek itu memenuhi
syarat. Objek yang ditunjuk disebut \emph{saksi} (\emph{witness}).

\begin{teorema}
Untuk setiap bilangan asli $n$ terdapat $n$ bilangan bulat berurutan yang semuanya
bukan prima.
\end{teorema}
\par\smallskip\noindent\textit{Bukti.} Ambil $N = (n+1)!$ dan perhatikan $n$ bilangan
\[
  N+2,\; N+3,\; \dots,\; N+(n+1).
\]
Untuk setiap $k$ dengan $2 \leq k \leq n+1$, bilangan $k$ membagi $(n+1)!$ karena $k$
adalah salah satu faktornya, sehingga $k$ juga membagi $N+k$. Selain itu
$N+k > k > 1$, jadi $N+k$ memiliki pembagi positif $k$ yang bukan $1$ dan bukan
dirinya sendiri. Menurut definisi, $N+k$ bukan prima. Kita telah menunjuk $n$ bilangan
yang diminta, dan tidak perlu mencari yang terkecil.
\hfill$\blacksquare${\parfillskip=0pt\par}

Untuk $n = 5$ bukti itu menunjuk $6! = 720$, lalu $722, 723, 724, 725, 726$. Kelimanya
memang terurai: $722 = 2 \cdot 19^{2}$, $723 = 3 \cdot 241$, $724 = 2^{2} \cdot 181$,
$725 = 5^{2} \cdot 29$, dan $726 = 2 \cdot 3 \cdot 11^{2}$. Bilangan $721 = 7 \cdot 103$ juga terurai, tetapi ia tidak termasuk daftar kita.

Bukti konstruktif tidak menjanjikan saksi terkecil atau tercepat; yang dinilai hanya
apakah objek yang ditunjuk memenuhi syarat. Menyangkal pernyataan berbentuk $\exists$
menuntut ditunjukkannya bahwa tidak satu pun objek memenuhi, dan di situlah kontradiksi
dibutuhkan.

\begin{resep}
Empat langkah baku membuktikan ``jika $P$ maka $Q$''.
\begin{enumerate}
  \item Tulis $P$ dan $Q$ menurut definisinya, masing-masing satu baris.
  \item Buka dengan ``andaikan $P$ benar''. Yang boleh diandaikan hanya $P$; $Q$ belum
        boleh dipakai.
  \item Turunkan $Q$ dengan substitusi dan aljabar, satu langkah per baris, dan sebut
        definisi yang dipakai.
  \item Tutup dengan kalimat yang menyebut $Q$ persis seperti bentuknya di langkah 1.
\end{enumerate}
Peringatan urutan: langkah 2 tidak boleh ditukar dengan langkah 3. Coretan dari kesimpulan
ke belakang boleh, tetapi bukti yang diserahkan harus berjalan dari $P$ ke $Q$; kalau
tidak, yang terbukti adalah konversnya.
\end{resep}

\section{Induksi dan induksi kuat}
\tingkat{menengah}

Induksi adalah pola bukti untuk pernyataan yang berderet, satu untuk setiap bilangan
asli. Kita menegakkan kartu pertama dan memastikan setiap kartu yang jatuh menjatuhkan
kartu berikutnya; setelah keduanya, tidak ada kartu yang boleh berdiri.

\kamus{induksi}{induction}{seperti deretan domino: cukup pastikan kartu pertama jatuh
dan setiap kartu menjatuhkan kartu di depannya, lalu seluruh deret jatuh tanpa perlu
didorong satu per satu.}

Prinsip induksi kita ambil sebagai sifat dasar $\mathbb{N} = \{1,2,3,\dots\}$, setara
dengan pernyataan bahwa setiap himpunan bagian tak kosong dari $\mathbb{N}$ mempunyai
anggota terkecil; tidak ada yang dibuktikan secara melingkar, kita hanya memilih titik
berangkat seperti pada Bab~3.

\begin{catatan}[Prinsip induksi]
Misalkan $P(n)$ pernyataan tentang bilangan asli. Jika $P(1)$ benar dan, untuk setiap
$k \in \mathbb{N}$, kebenaran $P(k)$ memaksa kebenaran $P(k+1)$, maka $P(n)$ benar
untuk setiap $n \in \mathbb{N}$.
\end{catatan}

Dengan prinsip itu, kita buktikan rumus yang sudah ditebak pada pembuka bagian ini.

\begin{teorema}
Untuk setiap bilangan asli $n$ berlaku
\[
  1 + 3 + 5 + \cdots + (2n-1) = n^{2}.
\]
\end{teorema}
\par\smallskip\noindent\textit{Bukti.} Basis: untuk $n = 1$ ruas kiri bernilai $1$ dan
ruas kanan bernilai $1^{2} = 1$. Langkah: andaikan rumus benar untuk $n = k$, yaitu
$1+3+\cdots+(2k-1) = k^{2}$. Suku berikutnya adalah $2(k+1)-1 = 2k+1$, sehingga
\[
  1+3+\cdots+(2k-1)+(2k+1) = k^{2} + 2k + 1 = (k+1)^{2}.
\]
Jadi rumus benar untuk $k+1$. Menurut prinsip induksi, rumus benar untuk setiap
$n \in \mathbb{N}$.
\hfill$\blacksquare${\parfillskip=0pt\par}

Contoh kedua berbentuk pertidaksamaan. Dari $2^{k} \geq k+1$ dan syarat $k \geq 0$
diperoleh $2^{k+1} = 2 \cdot 2^{k} \geq 2(k+1) \geq k+2$, sehingga dengan basis $n = 0$
berlaku $2^{n} \geq n+1$ untuk setiap $n \geq 0$; hipotesis dan syarat samping dipakai
sekaligus.

Induksi biasa menuntut satu langkah mundur. Kadang kesimpulan untuk $n$ hanya turun dari
banyak nilai sebelumnya sekaligus, dan di situ induksi biasa tidak cukup.

\begin{catatan}[Prinsip induksi kuat]
Jika untuk setiap $n \in \mathbb{N}$ kebenaran $P(1), P(2), \dots, P(n-1)$ memaksa
kebenaran $P(n)$, maka $P(n)$ benar untuk setiap $n \in \mathbb{N}$.
\end{catatan}

\kamus{induksi kuat}{strong induction}{seperti menyusun menara dari balok: untuk
meletakkan balok kesepuluh, kita boleh memakai seluruh balok di bawahnya, bukan hanya
balok kesembilan.}

Dengan izin itu kita dapat membuktikan hal yang tidak terjangkau induksi biasa.

\begin{teorema}
Setiap bilangan bulat $n \geq 2$ dapat ditulis sebagai hasil kali bilangan-bilangan
prima.
\end{teorema}
\par\smallskip\noindent\textit{Bukti.} Kita pakai induksi kuat pada $n$. Untuk $n = 2$:
bilangan $2$ sendiri prima, jadi ia hasil kali dengan satu faktor. Sekarang ambil
$n \geq 3$ dan andaikan pernyataan benar untuk semua bilangan bulat $m$ dengan
$2 \leq m < n$. Jika $n$ prima, selesai. Jika $n$ bukan prima, menurut definisi terdapat
pembagi positif $a$ dengan $a \neq 1$ dan $a \neq n$; tulis $n = ab$ dengan
$b = n/a$. Karena $a \neq 1$ dan $a \neq n$, berlaku $2 \leq a \leq n-1$, dan begitu
pula $2 \leq b \leq n-1$. Hipotesis induksi berlaku untuk $a$ maupun $b$, jadi keduanya
hasil kali bilangan prima, dan hasil kali keduanya adalah $n$.
\hfill$\blacksquare${\parfillskip=0pt\par}

Hipotesis induksi kuat dipakai dua kali, pada $a$ dan pada $b$. Di situlah induksi biasa
tidak nyaman: keduanya bisa jauh lebih kecil daripada $n-1$, sehingga pernyataan untuk
$n-1$ saja tidak menutupnya.

Teorema itu menutup janji bagian sebelumnya: daftar prima tidak pernah lengkap.

\begin{akibat}[Tak berhingga banyak bilangan prima]
\index{prima tak berhingga}
Untuk setiap daftar berhingga bilangan prima $p_{1}, p_{2}, \dots, p_{k}$ terdapat
bilangan prima yang tidak ada di dalam daftar itu.
\end{akibat}
\par\smallskip\noindent\textit{Bukti.} Misalkan $N = p_{1}p_{2}\cdots p_{k} + 1$. Karena
$N \geq 2$, teorema sebelumnya memberi bahwa $N$ adalah hasil kali bilangan prima;
ambil satu faktor prima $q$ dari $N$. Andaikan $q$ termasuk daftar, misalnya
$q = p_{i}$. Maka $q$ membagi hasil kali $p_{1}\cdots p_{k}$, dan karena $q$ juga
membagi $N$, bilangan $q$ membagi selisih $N - p_{1}\cdots p_{k} = 1$. Tidak ada
bilangan prima yang membagi $1$, jadi pengandaian itu salah. Maka $q$ tidak ada di
dalam daftar.
\hfill$\blacksquare${\parfillskip=0pt\par}

Bukti itu tidak membuktikan bahwa $N$ prima, hanya bahwa $N$ punya faktor prima di luar
daftar: untuk $2,3,5,7$ diperoleh $N = 211$ yang kebetulan prima, untuk $2,3,5,7,11,13$
diperoleh $N = 30031 = 59 \cdot 509$. Bukti itu konstruktif; ``andaikan hanya ada
berhingga banyak prima'' hanyalah pembungkus retorisnya.

\begin{untukanak}
Bayangkan deretan seribu kotak, dan kita ingin membuktikan setiap kotak berisi
kelereng. Kita tidak perlu membuka seribu kotak. Cukup buka kotak pertama, lalu buktikan
satu hal: kalau sebuah kotak berisi kelereng, kotak di sebelahnya juga berisi kelereng.
Kalau dua hal itu benar, seribu kotak selesai diperiksa tanpa dibuka satu per satu.

Inilah induksi, dan ia terasa seperti sulap. Kuncinya pada janji kedua: kalau hubungan
antarkotak bocor di satu tempat saja, seluruh deretan roboh.

Sekarang bayangkan kotak pertama ternyata kosong. Hubungan antarkotak tetap benar, tetapi
tidak menolong, karena tidak ada yang pertama untuk dijatuhkan. Inilah kesalahan yang
paling sering terjadi.

Di bab ini kita akan tahu mengapa dua langkah itu tidak boleh ditukar, dan mengapa
pemeriksaan pada banyak kotak bukan pengganti keduanya.
\end{untukanak}

\section{Ketika bukti gagal}
\tingkat{lanjut}

Kembali ke papan tulis pada pembuka bab. Empat puluh nilai pertama prima, dan pada
$n = 40$ pola itu patah:
\[
  40^{2} + 40 + 41 = 1600 + 40 + 41 = 1681 = 41^{2}.
\]
Kegagalan itu bukan kebetulan, dan bentuknya dapat dijelaskan tanpa komputer. Untuk
setiap bilangan bulat $p \geq 2$, substitusi $n = p-1$ memberi
\[
  (p-1)^{2} + (p-1) + p = p^{2} - 2p + 1 + p - 1 + p = p^{2},
\]
sedangkan substitusi $n = p$ memberi $p^{2} + p + p = p(p+2)$. Jadi setiap bilangan
$p \geq 2$ menghasilkan dua titik gagal, yaitu $n = p-1$ dan $n = p$. Dengan $p = 41$
keduanya menghasilkan $1681 = 41^{2}$ dan $1763 = 41 \cdot 43$.

Pola yang sama muncul pada polinomial yang lebih kecil. Untuk $n^{2}+n+11$ titik gagal
pertama $n = 10$, yang memberi $121 = 11^{2}$; untuk $n^{2}+n+17$ titik gagal pertama
$n = 16$, yang memberi $289 = 17^{2}$; untuk $n^{2}+n+29$ kegagalan datang pada $n = 2$,
karena $35 = 5 \cdot 7$ bukan prima. Semakin besar bilangan pada posisi konstanta,
semakin panjang rentang awal yang bersih; itulah seluruh rahasia reputasi $41$.

Kegagalan itu memberi tiga pelajaran yang sering tertukar.

Pelajaran pertama bersifat asimetris. Satu contoh penyangkal menumbangkan pernyataan
``untuk semua'', sedangkan empat puluh contoh yang berhasil tidak membuktikan apa pun
tentang nilai keempat puluh satu. Pemeriksaan komputer karena itu alat menemukan dugaan,
bukan alat membuktikan.

Pelajaran kedua menyangkut letak cacat. Cacat pertama yang paling sering lolos muncul
pada bukti induksi. Perhatikan klaim yang salah berikut: untuk setiap $n \geq 1$
berlaku $1+3+\cdots+(2n-1) = n^{2}+1$. Langkah induksinya, bila diperiksa, ternyata
sah:
\[
  \bigl(n^{2}+1\bigr) + (2n+1) = n^{2} + 2n + 2 = (n+1)^{2} + 1 .
\]
Jadi dari kebenaran rumus pada $n$ memang mengalir kebenaran pada $n+1$. Yang gagal
adalah basis: untuk $n = 1$ ruas kiri bernilai $1$ sedangkan ruas kanan bernilai $2$.
Langkah induksi memindahkan kebenaran, ia tidak menciptakannya.

Pelajaran ketiga menyangkut batas pemeriksaan. Uji pembagian lazim: $m > 1$ prima bila
tidak habis dibagi bilangan prima mana pun yang lebih kecil daripada $\sqrt{m}$. Untuk $m = 1681$ diperoleh $\sqrt{1681} = 41$, dan prima yang
lebih kecil daripada $41$ adalah $2, 3, 5, 7, 11, 13, 17, 19, 23, 29, 31, 37$; tidak
satu pun membagi $1681$. Dengan aturan itu $1681$ lolos sebagai prima. Kesalahannya ada
pada kata ``lebih kecil'': aturan yang benar memakai ``lebih kecil atau sama dengan'',
karena $1681 = 41 \cdot 41$ dan pembaginya tepat sebesar akar. Satu tanda pertidaksamaan
yang longgar menghasilkan satu prima palsu, tepat di $n = 40$.

Tiga cacat lain berulang di tulisan yang lebih panjang. Cacat penalaran melingkar:
``karena $n^{2}+n+41$ prima untuk semua $n$, maka prima juga untuk $n = 40$'' memakai
yang harus dibuktikan sebagai alasan. Cacat membagi dengan nol: dari $a = b$ orang
menyimpulkan $a+b = b$ setelah membagi dengan $a-b$, padahal $a-b = 0$. Cacat menukar
arah: dari $P \Rightarrow Q$ dan $Q$ orang menyimpulkan $P$; pada tabel kebenaran,
baris ketiga adalah tempat kesalahan itu bersembunyi.

Keempat cacat itu berbagi satu akar: sebuah langkah dipakai di luar syaratnya.
Kebiasaan yang berguna bukan menghafal pola, melainkan memeriksa syarat setiap langkah:
pada kontradiksi pengandaiannya harus negasi kesimpulan, pada induksi ada basis dan
langkah, pada pembagian pembaginya tidak nol.

\begin{tantangan}
Klaim: untuk setiap bilangan bulat $k \geq 0$, bilangan $2^{2^{k}}+1$ prima. Untuk
$k = 0,1,2,3,4$ nilainya $3$, $5$, $17$, $257$, $65537$, dan kelimanya memang prima.
Tentukan $k$ terkecil yang menggugurkan klaim itu, lalu tunjukkan pembaginya; klaim itu
salah, dan kesalahannya dapat ditunjukkan tanpa komputer.

Petunjuk 1: lima nilai bukan bukti. Periksa $k = 5$; bilangannya $4294967297$, dan pembagi
terkecilnya jauh lebih kecil daripada akar kuadratnya.

Petunjuk 2: pembagi itu berbentuk $5 \cdot 2^{7}+1 = 641$, dan
$4294967297 = 641 \cdot 6700417$.

Kalau mentok, kembali ke bagian ``Ketika bukti gagal''.
\end{tantangan}

\section{Contoh hitung}
\tingkat{menengah}

Tiga perhitungan berikut menopang bab ini dan dapat diulang di kertas. Yang pertama
jumlah bilangan ganjil pertama, suku demi suku dari $n = 1$.

{\centering
\begin{tabular}{cccc}
\hline
$n$ & suku-suku & jumlah & $n^{2}$ \\
\hline
$1$ & $1$ & $1$ & $1$ \\
$2$ & $1+3$ & $4$ & $4$ \\
$3$ & $1+3+5$ & $9$ & $9$ \\
$4$ & $1+3+5+7$ & $16$ & $16$ \\
$5$ & $1+3+5+7+9$ & $25$ & $25$ \\
$6$ & $1+3+5+7+9+11$ & $36$ & $36$ \\
\hline
\end{tabular}\par}

Lanjutannya: $36+13 = 49 = 7^{2}$, lalu $49+15 = 64$, $64+17 = 81$, $81+19 = 100$.
Selisih dua kuadrat berurutan memang $(n+1)^{2}-n^{2} = 2n+1$, dan di situlah induksi
berpijak.

Yang kedua uji primalitas pada $n^{2}+n+41$, hanya pada nilai yang menarik.

{\centering
\begin{tabular}{ccc}
\hline
$n$ & $n^{2}+n+41$ & keterangan \\
\hline
$0$ & $41$ & prima \\
$1$ & $43$ & prima \\
$39$ & $1601$ & prima \\
$40$ & $1681 = 41^{2}$ & bukan prima \\
$41$ & $1763 = 41 \cdot 43$ & bukan prima \\
\hline
\end{tabular}\par}

Pemeriksaan menyeluruh untuk $n = 0$ sampai $n = 40$ memberi $40$ nilai prima dari $41$;
satu-satunya yang gagal adalah $n = 40$, dan $\sqrt{1681} = 41$ tepat, sehingga pembagi
yang menentukan berada di batas pemeriksaan.

Yang ketiga pemeriksaan kontraposisi pada tabel kebenaran, ditulis lengkap.

{\centering
\begin{tabular}{cc|cc|c}
\hline
$P$ & $Q$ & $P \Rightarrow Q$ & $\neg Q \Rightarrow \neg P$ & setara \\
\hline
B & B & B & B & ya \\
B & S & S & S & ya \\
S & B & B & B & ya \\
S & S & B & B & ya \\
\hline
\end{tabular}\par}

Keempat baris memberi kesamaan: pola kontraposisi bukan trik retoris melainkan penggantian
bentuk yang tidak mengubah nilai. Kolom $Q \Rightarrow P$ akan berbeda pada baris ketiga,
dan perbedaan itulah yang membuat konvers tidak boleh dipakai.

\kertas{Hitung $1+3+\cdots+(2n-1)$ untuk $n = 11$ dan $n = 12$ dengan bertolak dari
$100 = 10^{2}$, yaitu tambahkan dua bilangan ganjil berikutnya satu per satu. Lalu
hitung $n^{2}+n+17$ untuk $n = 16$ dan periksa dengan pembagian apakah $289$ sama
dengan $17^{2}$. Terakhir, cari titik gagal pertama $n^{2}+n+11$ dengan menguji
$n = 0, 1, 2, \dots$ dan jelaskan mengapa titik itu dapat diramalkan dari bentuk
polinomialnya.}

\section{Lab komputasi}
\tingkat{lanjut}

Skrip berikut hanya memakai pustaka standar Python dan dapat dijalankan ulang tanpa
pemasangan apa pun; setiap baris keluarannya adalah keluaran asli.

\begin{lab}
\begin{verbatim}
from itertools import product

def prima(n):
    if n < 2:
        return False
    d = 2
    while d * d <= n:
        if n % d == 0:
            return False
        d += 1
    return True

pasang = product([False, True], repeat=2)
print("kontraposisi:", all((not P or Q) == (Q or not P)
                          for P, Q in pasang))

D = range(1, 11)
print("y>x:", sum(y > x for x in D for y in D), "dari 100")
print("x tanpa saksi:", [x for x in D if all(y <= x for y in D)])

nilai = [(n, n*n + n + 41) for n in range(41)]
print("prima:", sum(prima(v) for _, v in nilai), "dari 41")
print("gagal:", [(n, v, int(v**0.5)) for n, v in nilai if not prima(v)])
print("41*41, 41*43:", 41*41, 41*43)

sisa = {(n*n - 1) % 8 for n in range(1, 40, 2)}
print("sisa mod 8:", sisa)

langkah = all(n*n + 1 + 2*n + 1 == (n+1)**2 + 1 for n in range(1, 500))
print("langkah palsu:", langkah)
\end{verbatim}
\end{lab}

Keluaran yang diperoleh:

\begin{verbatim}
kontraposisi: True
y>x: 45 dari 100
x tanpa saksi: [10]
prima: 40 dari 41
gagal: [(40, 1681, 41)]
41*41, 41*43: 1681 1763
sisa mod 8: {0}
langkah palsu: True
\end{verbatim}

Baris pertama memeriksa keempat baris tabel sekaligus. Baris kedua dan ketiga menghitung
yang tidak nyaman dikerjakan di kertas: dari seratus pasangan $(x,y)$ pada
$\{1,\dots,10\}$, $45$ memenuhi $y > x$, dan satu-satunya $x$ tanpa saksi adalah $10$;
itulah kegagalan $(\forall x)(\exists y)$ pada himpunan berhingga.

Baris keempat dan kelima memberi $40$ prima dari $41$ nilai, gagal hanya di
$(40,1681,41)$; angka $41$ itu adalah $\sqrt{1681}$ yang dipotong. Baris keenam
memastikan $1681 = 41^{2}$ dan $1763 = 41 \cdot 43$. Baris ketujuh menyatakan seluruh
$n$ ganjil di bawah $40$ memberi sisa $0$ modulo $8$, sesuai teorema pola pertama. Baris
terakhir peringatan: langkah induksi klaim yang salah itu sah untuk lima ratus nilai
pertama, sehingga pemeriksaan numerik tidak menemukan cacatnya.

\begin{untukdosen}
Urutan mengajar yang berhasil dipakai: mulai dari $n = 40$ dan $1681$, bukan dari
definisi. Setelah mahasiswa kesal karena empat puluh pemeriksaan ternyata sia-sia, baru
masuk tabel kebenaran, dengan penekanan pada baris ketiga dan keempat. Lanjutkan ke
negasi kuantor lewat definisi limit barisan, lalu ke definisi harfiah dengan contoh
$91 = 7 \cdot 13$; tutup dengan induksi, kegagalan basis, dan lab.

Dua miskonsepsi tersering. Pertama, implikasi dibaca sebagai sebab-akibat, sehingga baris
dengan anteseden salah dinilai salah dan bukti berkasus hampa ditolak. Kedua, induksi
dianggap selesai setelah langkah terbukti, dan induksi kuat dianggap induksi biasa yang
diulang; akibatnya basis tidak diperiksa dan hipotesis dipakai tanpa menyebut
rentangnya.

Aktivitas papan tulis: bagi kelas menjadi dua kelompok. Kelompok pertama menulis klaim
berkuantor, kelompok kedua menulis negasinya, lalu masing-masing menyebut satu objek yang
menyangkal klaim lawan. Permainan berhenti ketika sebuah negasi tidak dapat disangkal.

Mahasiswa siap lanjut apabila dapat menegasikan definisi limit tanpa membuka catatan
dan dapat menjelaskan mengapa basis pada induksi tidak boleh dilewati.
\end{untukdosen}

\section{Ringkasan}
\tingkat{dasar}

\begin{ringkasan}
\begin{itemize}
  \item Pernyataan bernilai benar atau salah; $P \Rightarrow Q$ hanya salah ketika $P$
        benar dan $Q$ salah.
  \item $P \Rightarrow Q \equiv \neg Q \Rightarrow \neg P$, tetapi konversnya tidak
        setara; $\neg(P \Rightarrow Q) \equiv P \wedge \neg Q$.
  \item Negasi berjalan dari luar ke dalam: setiap kuantor berganti jenis, sehingga
        menyangkal ``untuk semua'' berarti menunjuk satu contoh penyangkal.
  \item Definisi berlaku dua arah: nama hanya boleh dipakai bila seluruh klausa
        terpenuhi, dan seluruh klausa dipakai untuk memeriksanya.
  \item Empat pola bukti: langsung, kontraposisi, kontradiksi, konstruktif. Pilih
        yang paling pendek, bukan yang paling megah.
  \item Induksi menuntut basis dan langkah; induksi kuat menuntut seluruh nilai
        sebelumnya. Langkah tanpa basis tidak membuktikan apa pun.
\end{itemize}
\end{ringkasan}

\section{Latihan}
\tingkat{menengah}

\begin{latihan}
\begin{enumerate}
  \item Tulis negasi dari pernyataan berikut, lalu terjemahkan hasilnya menjadi
        kalimat Indonesia yang menyatakan apa yang harus ditunjukkan untuk
        menyanggahnya: ``untuk setiap $\varepsilon > 0$ terdapat bilangan asli $N$
        sehingga untuk setiap $n \geq N$ berlaku $|a_{n} - L| < \varepsilon$''.

  \item Tulis negasi dari ``setiap bilangan prima lebih besar daripada $2$ adalah
        ganjil'', lalu sebutkan satu objek yang menyangkal pernyataan aslinya.
        Kerjakan hal yang sama untuk ``terdapat bilangan asli $n$ sehingga
        $n^{2}+n+41$ tidak prima''.

  \item Buktikan bahwa untuk setiap bilangan bulat $n$, bilangan $n^{2}+n$ selalu
        genap. Kerjakan dengan dua cara: pembagian kasus paritas, dan kontraposisi.

  \item Buktikan dengan induksi bahwa untuk setiap bilangan asli $n$ berlaku
        \[
          1^{2} + 2^{2} + \cdots + n^{2} = \frac{n(n+1)(2n+1)}{6}.
        \]
        Periksa rumus itu untuk $n = 1, 2, 3$ sebelum menyusun buktinya.

  \item Buktikan dengan induksi kuat bahwa setiap bilangan asli $n$ dapat ditulis
        dalam bentuk $n = 2^{k}m$ dengan $k \geq 0$ bilangan bulat dan $m$ bilangan
        ganjil. Tunjukkan di bagian mana hipotesis induksi kuat dipakai.

  \item Nilai klaim berikut: benar, salah, atau belum cukup bukti.
        \begin{enumerate}
          \item Jika $P(1)$ benar dan $P(k) \Rightarrow P(k+1)$ untuk setiap
                $k \in \mathbb{N}$, maka $P(n)$ benar untuk setiap $n \in
                \mathbb{N}$.
          \item Jika $P(n)$ benar untuk $n = 1, 2, \dots, 10^{6}$, maka $P(n)$ benar
                untuk setiap $n \in \mathbb{N}$.
          \item Jika $n^{2}+n+41$ prima untuk $n = 0, 1, \dots, 39$, maka klaim itu
                benar untuk setiap $n \geq 0$.
          \item Jika $3$ membagi $n$, maka $6$ membagi $n$.
        \end{enumerate}
        Untuk setiap butir yang salah, berikan satu contoh penyangkal konkret.

  \item Hitung dengan Python banyaknya bilangan prima di antara $n^{2}+n+41$ untuk
        $n = 0$ sampai $n = 100$, lalu tentukan semua $n$ pada rentang itu yang
        membuat nilainya bukan prima. Periksa apakah setiap titik gagal itu berbentuk
        $n = p-1$ atau $n = p$ untuk suatu bilangan prima $p$.

  \item Untuk $n = 1$ sampai $n = 12$, hitung $2^{n}$ dan $n^{2}$, lalu tentukan
        bilangan asli terkecil $n_{0}$ sehingga $2^{n} > n^{2}$ berlaku untuk setiap
        $n \geq n_{0}$. Buktikan klaim itu dengan induksi, dengan basis $n_{0}$.

  \item Buktikan bahwa pernyataan ``terdapat bilangan bulat $n > 1$ dengan
        $n^{2}-1$ prima'' salah. Sebutkan lebih dulu bentuk faktorisasi $n^{2}-1$
        untuk $n > 1$, lalu jelaskan mengapa faktor-faktornya tidak mungkin sama
        dengan $1$.

  \item Tunjukkan bahwa pernyataan ``jika $n^{2}$ habis dibagi $4$ maka $n$ habis
        dibagi $4$'' salah, lalu perbaiki kesimpulannya menjadi klaim yang benar dan
        buktikan.
\end{enumerate}

\par\smallskip\noindent\textit{Petunjuk untuk tiga butir tersulit.} Butir 5: pisahkan
kasus $n$ ganjil, yang langsung selesai dengan $k=0$ dan $m=n$, dari kasus $n$ genap,
yang memakai hipotesis induksi pada $n/2$. Butir 7: hitung juga $\sqrt{n^{2}+n+41}$
untuk setiap titik gagal, dan bandingkan dengan pembagi terkecilnya. Butir 9: $n^{2}-1 = (n-1)(n+1)$; untuk $n > 1$ kedua faktor itu terletak di antara $1$
dan $n^{2}-1$.
\end{latihan}

\begin{untukpeneliti}
Yang benar-benar terbukti di bab ini tiga hal klasik yang bertahan tanpa komputer:
ketakrasionalan $\sqrt{2}$ dengan bukti paritas, tak berhingga banyak prima lewat
konstruksi Euklides, dan ketunggalan bentuk setiap bilangan asli sebagai hasil kali prima.
Yang hanya dipopulerkan adalah cerita bahwa $41$ ``hampir'' rumus prima; yang terjadi
hanyalah rentang awal yang panjang, dan panjangnya dapat diramalkan dari $n^{2}+n+p$.

Sumber primer: Abbott, \emph{Understanding Analysis}, Bab~1 Bagian~1.1 (irasionalitas
$\sqrt{2}$) dan Bagian~1.2 (pendahuluan, termasuk induksi); Avigad dan Massot,
\emph{Mathematics in Lean}, Bab~3 (logika) serta Bab~5 Bagian~5.1 dan 5.2
(akar irasional, induksi dan rekursi); Euklides, \emph{Elements}, Buku~IX
Proposisi~20 untuk tak berhingga prima.

Pertanyaan yang masih terbuka: tidak ada polinomial berderajat positif dengan koefisien
bulat yang bernilai prima untuk semua $n \geq 0$ (bukti satu baris: $f(0)$ membagi
$f(t\,f(0))$), tetapi seberapa panjang rentang prima yang dapat dicapai polinomial
berderajat tinggi masih dipelajari, dan distribusi prima terikat pada hipotesis Riemann
yang belum terbukti.

Batas keberlakuan: semua hasil bab ini elementer dan tidak bergantung pada pilihan aksioma
himpunan. Verifikasi formal di Lean (Bab~28) memeriksa keabsahan setiap langkah deduktif,
tetapi tidak memeriksa apakah pernyataan formal yang dibuktikan memang yang dimaksud
penulisnya.
\end{untukpeneliti}

\peta{Bab 1 pada \emph{Understanding Analysis} (Abbott), khususnya Bagian 1.1 tentang
irasionalitas $\sqrt{2}$ dan Bagian 1.2 tentang pendahuluan; Bab 3 (logika), Bab 5
Bagian 5.1 dan 5.2 (akar irasional serta induksi dan rekursi) pada
\emph{Mathematics in Lean} (Avigad dan Massot); dan Buku IX Proposisi 20
\emph{Elements} (Euklides) untuk tak berhingga banyak bilangan prima.}
